\section*{Hoofdstuk \ref{ch:b3:plant-molecular-physiology} --- Moleculaire plantenfysiologie en stressreacties}

\begin{solution}{exo:b3:plant-molecular-physiology:1}
Fytochromen, rood \qty{660}{nm} en verrood \qty{730}{nm}:
\omterm{def:b3:plant-molecular-physiology:photoreceptors}{de-etiolering}, \omterm{prop:b3:plant-molecular-physiology:photoequilibrium}{schaduwontwijking}, kieming. Cryptochromen, blauw
\qty{450}{nm}: \omterm{def:b3:plant-molecular-physiology:photoreceptors}{de-etiolering}, de klok, de bloei. \omterm{def:b3:plant-molecular-physiology:photoreceptors}{Fototropinen}, blauw:
buigen naar het licht, openen van de huidmondjes, beweging van de
chloroplasten. Chlorofyl meet licht als energie en doet niets met de
informatie; de fotoreceptoren lezen kleur, richting en duur bij stromen
die een miljoen maal lager zijn en veranderen de genexpressie --- het
verschil tussen het licht opeten en het lezen.
\end{solution}

\begin{solution}{exo:b3:plant-molecular-physiology:2}
Elk van deze \omterm{def:b3:endocrinology:hormone}{hormonen} bevordert de ubiquitinering en vernietiging van
een repressor: \omterm{def:b3:plant-molecular-physiology:hormones}{auxine} lijmt \omterm{def:b3:plant-molecular-physiology:hormones}{TIR1} aan de Aux/IAA-eiwitten, \omterm{def:b3:plant-molecular-physiology:hormones}{gibberelline}
gebonden aan GID1 levert de DELLA-eiwitten aan een F-box-ligase af,
jasmonaat lijmt COI1 aan de JAZ-eiwitten. De repressoren zitten al op
hun doelgenen met de activatoren ernaast klaar, dus hoeft er geen nieuw
eiwit te worden gemaakt: de repressor vernietigen, wat het proteasoom
in minuten doet, zet de genen meteen aan.
\end{solution}

\begin{solution}{exo:b3:plant-molecular-physiology:3}
ABA bindt PYR/PYL; het complex remt het fosfatase PP2C; het kinase
SnRK2 (OST1), dat niet meer wordt gedefosforyleerd, wordt actief en
fosforyleert het anionkanaal SLAC1; anionen vertrekken, het membraan
depolariseert, uitwaartse kaliumkanalen gaan open, K$^{+}$ en daarna
water vertrekken, de turgor daalt en de \omterm{prop:b3:plant-molecular-physiology:stomata}{sluitcellen} zakken tegen elkaar
aan. Elke schakel kan breken: geen receptoren, een fosfatase dat niet
kan worden geremd (de klassieke mutant \emph{abi1}), geen OST1 of geen
SLAC1 --- elk geeft een slappe plant wier huidmondjes bij droogte open
blijven staan.
\end{solution}

\begin{solution}{exo:b3:plant-molecular-physiology:4}
PTI herkent moleculen die hele klassen microben gemeen hebben
(flagelline, chitine) met receptorkinasen aan het celoppervlak en
eindigt in een versterkte wand, gesloten huidmondjes, expressie van
afweergenen en tragere groei, gewoonlijk zonder celdood. ETI herkent
een bepaald effectoreiwit of de schade ervan met een intracellulaire
\omterm{def:b3:plant-molecular-physiology:immunity}{NLR-receptor} en eindigt in de overgevoelige dood van de besmette cellen
en een alarm door de hele plant.
\end{solution}

\begin{solution}{exo:b3:plant-molecular-physiology:5}
Daglicht $0.87\times 1.15/1.75 = 0.57$; één blad $0.87\times 0.4/1.0 =
0.35$; dicht bladerdak $0.87\times 0.1/0.7 = 0.12$; gloeilamp
$0.87\times 0.7/1.3 = 0.47$. Een gloeidraadlamp is rijk aan verrood,
dus ligt $\varphi$ onder die van het daglicht en leest de kiemplant
gedeeltelijke schaduw: zij strekt zich. Koelwitte lampen en leds zonder
verrood doen het omgekeerde.
\end{solution}

\begin{solution}{exo:b3:plant-molecular-physiology:6}
Wand $5.0$, cytosol $7.5$: $(1 + 10^{2.75})/(1 + 10^{0.25}) = 563/2.78
= 203$. Wand $4.5$: $563/(1 + 10^{-0.25}) = 563/1.56 = 360$ --- de wand
aanzuren verhoogt het neutrale deel buiten en daarmee de opname.
Cytosol $6.5$ bij een wand van $5.5$: $(1 + 10^{1.75})/6.62 = 57/6.62 = 8.6$:
de val wordt vijf maal zwakker, \omterm{def:b3:plant-molecular-physiology:hormones}{auxine} lekt langs elk vlak weg en de
polaire vervallen vlakken af --- een reden dat wortels zonder zuurstof
hun oriëntatie verliezen.
\end{solution}

\begin{solution}{exo:b3:plant-molecular-physiology:7}
\qty{1}{cm/h} $= \qty{10000}{\micro m/h}$: $200$ cellen per uur,
\qty{18}{s} in elk. Diffusie over \qty{50}{\micro m}: $t =
(5\times 10^{-5})^{2}/(2\times 5\times 10^{-10}) = \qty{2.5}{s}$. Het
inwendige van de cel wordt in seconden doorkruist; het meeste van de
\qty{18}{s} gaat op aan het wachten om via de PIN's aan het basale
membraan naar buiten te worden gedragen, en dat is de
snelheidsbepalende stap.
\end{solution}

\begin{solution}{exo:b3:plant-molecular-physiology:8}
$E/A = 1.6\,\Delta w/\Delta c = 1.6\times\num{16000}/150 = 171$
watermoleculen per CO$_{2}$. C$_{4}$, $\Delta c = 250$: $102$. Droge
dag, $\Delta w = 24$: $256$.
\end{solution}

\begin{solution}{exo:b3:plant-molecular-physiology:9}
Vriezen droogt cellen uit doordat het water naar ijs buiten de cel
trekt, dus zijn koude en droogte allebei uitdrogingsstress. Zij delen
tweede boodschappers (ABA, calciumpieken, reactieve zuurstof),
transcriptiefactoren (de CBF/DREB-familie wordt door beide opgewekt) en
producten: dehydrinen en LEA-eiwitten die membranen en eiwitten tegen
waterverlies afschermen, verenigbare opgeloste stoffen (proline,
suikers) die water vasthouden, en antioxidanten. Een plant die ze voor
de ene stress heeft gemaakt heeft ze voor de andere.
\end{solution}

\begin{solution}{exo:b3:plant-molecular-physiology:10}
Bij \qty{1}{\%} van de volle zon is $k_{1} + k_{2} = \qty{0.005}{s^{-1}}$:
95~\% van het evenwicht na $3/k = \qty{600}{s}$, tien minuten, tegen
zes seconden op het middaguur. Acht uur is vier halveringstijden: $1/16 \approx
\qty{6}{\%}$ van het Pfr blijft over. Een korte flits zet
$\varphi$ wel, maar het Pfr vervalt daarna, terwijl een dag het uren
hoog houdt; reacties die vergen dat Pfr uren werkt (kieming,
\omterm{def:b3:plant-molecular-physiology:photoreceptors}{de-etiolering}) tellen de blootstelling op, zodat een zaad dat door een
ploeg even bloot komt te liggen niet antwoordt zoals een zaad dat aan
het oppervlak blijft liggen.
\end{solution}

\begin{solution}{exo:b3:plant-molecular-physiology:11}
Tien plekken per dag met $100$ cellen elk: \num{1000} cellen, $10^{-4}$
van het blad; over een seizoen een fractie van een procent, tegenover
het verlies van het blad als één besmetting zich ongehinderd
verspreidde. De dode cellen kosten bijna niets en nemen het voedsel van
de ziekteverwekker met zich mee --- voor een biotroof, die levende
cellen nodig heeft. Een necrotroof leeft van dood weefsel: de
\omterm{def:b3:plant-molecular-physiology:immunity}{overgevoeligheidsreactie} reikt hem een maaltijd en een steunpunt aan,
en sommige (\emph{Botrytis}, \emph{Sclerotinia}) scheiden gifstoffen
uit die haar met opzet uitlokken; tegen hen steunt de plant op de
afweer die door jasmonaat wordt geregeld en niet op celdood.
\end{solution}

\begin{solution}{exo:b3:plant-molecular-physiology:12}
Aanmaak met snelheid $p$ van uur 12 tot uur 16. Lange dag, de hele tijd
licht, $k = \ln 2/4 = \qty{0.17}{h^{-1}}$: de waarde op uur 16 is $(p/k)(1 -
\mathrm{e}^{-4k}) = (p/0.17)(0.5) = 2.9p$. Korte dag, donker vanaf uur
10, $k = \qty{1.39}{h^{-1}}$: $(p/1.39)(1 - \mathrm{e}^{-5.5})
\approx 0.72p$. Vier maal zoveel CONSTANS in de lange dag. Een drempel
daartussen --- zeg $2p$, nodig om \emph{FT} te activeren --- wordt
alleen overschreden wanneer het aanmaakvenster van de klok met licht
samenvalt: het samenvallen van een intern ritme met de uitwendige dag
zet een geleidelijke daglengte om in een alles-of-nietsbesluit tot
bloeien.
\end{solution}

\begin{solution}{pb:b3:plant-molecular-physiology:1}
\textbf{1.} Daglicht $0.87\times 1.15/1.75 = 0.57$; bladerdak $0.87\times
0.2/0.8 = 0.22$.
\textbf{2.} $r = 2(1 - \varphi)$: \qty{0.86}{mm/h} in daglicht,
\qty{1.57}{mm/h} in de schaduw; over \qty{48}{h}: $41$ tegen
\qty{75}{mm}, \qty{34}{mm} extra.
\textbf{3.} Het rood daalt tot $0.1\times 1.15 = 0.115$ van het
verrood, dat tot $0.8$ daalt: $\zeta = 0.144$, $\varphi = 0.87\times 0.144/0.744 =
0.17$. Ja: één blad laat $\varphi$ van $0.57$ naar $0.17$ zakken, diep
in het gebied van de \omterm{prop:b3:plant-molecular-physiology:photoequilibrium}{schaduwontwijking}.
\textbf{4.} Beide omzettingssnelheden schalen met de stroom, dus hun
verhouding niet. Onder bewolking is het spectrum vrijwel onveranderd en
blijft $\varphi$ rond $0.57$: een kiemplant in de open lucht op een
grauwe dag strekt zich niet --- zij leest buren, geen helderheid.
\textbf{5.} \qty{95}{\%} na $3/(k_{1} + k_{2})$: \qty{6}{s} in volle
zon, \qty{600}{s} bij \qty{1}{\%}.
\textbf{6.} \qty{8}{h} $=$ vier halveringstijden, $1/16 = \qty{6.3}{\%}$;
\qty{14}{h} $=$ zeven, $1/128 = \qty{0.8}{\%}$. Het Pfr dat bij
dageraad overblijft zou de lengte van de nacht kunnen melden, maar de
terugkeer is een thermische reactie die op warme nachten sneller loopt,
en de beginwaarde hangt af van het licht van de dag: planten meten de
nachtlengte in plaats daarvan met de klok.
\textbf{7.} Wand: $1/(1 + 10^{0.75}) = 0.15$; cytosol: $1/(1 +
10^{2.45}) = 0.0035$.
\textbf{8.} $(1 + 282)/(1 + 5.6) = 283/6.6 = 43$. Cytosol
$43\times 0.1 = \qty{4.3}{\micro M}$.
\textbf{9.} $\qty{1}{cm/h}/\qty{100}{\micro m} = 100$ cellen per uur,
\qty{36}{s} elk.
\textbf{10.} \qty{5}{h}. Het buigen begint binnen een uur omdat de
herverdeling die telt zijdelings is, over een millimeter top, en het
antwoord plaatselijk; de stroom over lange afstand bepaalt de aanvoer,
niet de timing.
\textbf{11.} Onderste flank: $60/50 = 1.2$, \omterm{def:b3:plant-molecular-physiology:hormones}{auxine} \qty{20}{\%} boven
de norm, strekking $0.8$; bovenste flank: $0.8$ van de norm, strekking
$1.2$. De bovenste flank groeit de onderste voorbij: de wortel buigt
omlaag.
\textbf{12.} Scheutcellen liggen onder hun optimum voor \omterm{def:b3:plant-molecular-physiology:hormones}{auxine}, dus
bevordert het extra \omterm{def:b3:plant-molecular-physiology:hormones}{auxine} op de onderste flank hun groei: de onderste
flank groeit de bovenste voorbij en de scheut buigt omhoog.
\textbf{13.} $E = 0.3\times 16 = \qty{4.8}{mmol\,m^{-2}\,s^{-1}}$.
\textbf{14.} $A = (0.3/1.6)\times 150 = \qty{28}{\micro mol\,m^{-2}\,s^{-1}}$;
$E/A = 4800/28 = 171$.
\textbf{15.} Oppervlak \qty{0.002}{m^2}, \num{43200}~s: water $4.8\times
10^{-3}\times 0.002\times\num{43200} = \qty{0.41}{mol} = \qty{7.5}{g}$;
CO$_{2}$ $28\times 10^{-6}\times 0.002\times\num{43200} =
\qty{2.4e-3}{mol}$, d.w.z.\ $2.4\times 10^{-3}/6\times 180 =
\qty{0.073}{g}$ glucose --- honderd gram water voor een gram suiker.
\textbf{16.} $E = 0.05\times 16 = \qty{0.8}{mmol\,m^{-2}\,s^{-1}}$,
$A = \qty{4.7}{\micro mol\,m^{-2}\,s^{-1}}$, verhouding nog altijd
$171$. De plant heeft zes zevende van haar water bespaard en zes
zevende van haar koolstof opgegeven: de klep sluiten verandert het
tempo, niet de prijs.
\textbf{17.} C$_{4}$: $\Delta c = 250$, $E/A = 4800/46.9 = 102$. CAM
's nachts: $E = 0.3\times 5 = \qty{1.5}{mmol}$, $A = 28$: $E/A = 53$.
\textbf{18.} $E/A = 1.6\,\Delta w/\Delta c$: de geleiding valt weg
omdat beide gassen door dezelfde opening gaan. De plant kan de
vervallen veranderen --- CO$_{2}$ binnen concentreren (C$_{4}$), openen
wanneer de lucht vochtig is (CAM, dageraad), het blad koelen, de
grenslaag met haren verdikken --- maar niet de natuurkunde van twee
gassen die één gat delen.
\textbf{19.} \num{1000} cellen per dag, $10^{-4}$ van het blad.
\textbf{20.} \num{60000} cellen in 60 dagen, \qty{0.6}{\%}. Eén ontsnapping
per dag die \qty{5}{\%} vernietigde zou het blad in twintig dagen
opsouperen: de \omterm{def:b3:plant-molecular-physiology:immunity}{overgevoeligheidsreactie} kost een honderdste van wat één
ontsnapte besmetting kost.
\textbf{21.} $0.2\times 5 = \qty{1}{\%}$ van de fotosynthese. Altijd
aangehouden zou zij \qty{5}{\%} kosten en zou de plant worden
voorbijgegroeid door buren die dat aan bladeren besteden; een opgewekte
afweer loont alleen wanneer de bedreiging er is.
\textbf{22.} Hij wint onzichtbaarheid voor die NLR en besmet het
resistente ras; hij verliest wat de effector deed (de PTI
onderdrukken), dus is hij zwakker op planten zonder de NLR. De
plantenpopulatie gaat NLR's tegen de overgebleven effectoren
bevoordelen, en het oude resistentiegen, nu nutteloos, loopt terug: een
frequentieafhankelijke cyclus van de genen van beide partijen.
\textbf{23.} De nabootsing (coronatine) activeert de jasmonaatroute,
die de salicylaatroute tegenwerkt; salicylaat is de afweer tegen
biotrofen zoals de bacterie zelf, en de nabootsing opent bovendien de
huidmondjes weer die de plant had gesloten. De ziekteverwekker keert de
ene arm van het afweersysteem tegen de andere.
\textbf{24.} Eén resistentiegen over miljoenen gelijke planten is een
eenvormige selectie: elke mutant die de effector verliest of verandert
verspreidt zich in een paar seizoenen over het hele gewas. Kwekers
stapelen verscheidene resistentiegenen zodat er meerdere effectoren
tegelijk verloren moeten gaan, mengen rassen, wisselen genen over de
jaren af, en bevoordelen de brede, door patronen opgewekte resistentie
waaraan geen enkele mutatie ontsnapt.
\textbf{25.} $\varphi \approx 0.22$ onder het bladerdak; \omterm{def:b3:plant-molecular-physiology:hormones}{auxine} binnen :
buiten $\approx 43$; ongeveer $170$ watermoleculen verloren per
vastgelegd CO$_{2}$.
\end{solution}
